% Quick start guide
\documentclass[aspectratio=11]{beamer}
% \usepackage[paperwidth=11.25n, paperheight=11.25in]{geometry}
% \geometry{total={11.25in,11.25in}}

% \usepackage[math]{iwona}
% \usepackage{cmbright}
% \usepackage{mathptmx}
% \usepackage[charter]{mathdesign}
% \usepackage{arev}
\RequirePackage{xcolor}


\usepackage{fontspec}
\setmonofont{Fira Code}[
  Contextuals=Alternate  % Activate the calt feature
]
\usepackage{listings}
\usepackage{lstfiracode} % https://ctan.org/pkg/lstfiracode


\usepackage{minted}

\usepackage{ulem}
\renewcommand{\ULdepth}{1.8pt}

\usepackage{bm}
\usepackage{textcomp}

\usepackage{dashrule}
% \usepackage{arydshln}
% \setlength\dashlinedash{0.2pt}
% \setlength\dashlinegap{1.5pt}
% \setlength\arrayrulewidth{0.3pt}

\usepackage{tikz}
\usetikzlibrary{snakes}
\usetikzlibrary{decorations.pathmorphing}
\usetikzlibrary{shadows}
\RequirePackage{tikz-penciline}


\usepackage[framemethod=TikZ]{mdframed}
\usepackage{showexpl}


\definecolor{NordBG}{HTML}{242933}

% set proper informations for pdf-file
\hypersetup{
	pdftitle=beamer-nord,
	pdfauthor=mfurquim,
	pdfsubject=presentation,
	hidelinks, % kill ugly link boxes
}

\usetheme{Nord}

% \renewcommand{\familydefault}{FiraCode Mono}
% \renewcommand{\familydefault}{\sfdefault}
% \renewcommand{\familydefault}{AuriocusKalligraphicus}


\title{\LaTeX\ Beamer introduction}
\subtitle{Quick-start guide}
\author{mateus@mfurquim.dev}
\institute{MFurquim Dev}
\date{\today}


% \logo{
% \begin{figure}
% 	\def\svgwidth{\columnwidth}
% 	\scalebox{0.5}{\input{logo.pdf_tex}}
% \end{figure}
% }

% kill navigation symbols in presentation
\beamertemplatenavigationsymbolsempty

\setbeamertemplate{headline}
{%
	% \vskip2pt\ \hskip2pt\tiny\textcolor{NordGreen}{\insertsection}\hfill\textcolor{NordBlue}{\insertsubsection}\hfill\ttfamily\insertpagenumber/\insertdocumentendpage\ \ \vskip2pt
	% \vskip2pt\ \hskip2pt\footnotesize\textcolor{NordGreen}{\insertsection}\hfill\textcolor{NordBlue}{\insertsubsection}\ \ \vskip2pt
    \vskip6pt\null\hfill\Large\textcolor{NordGreen}{\bfseries\insertsection}\hfill\null\newline%
    \vskip3pt\hfill\normalsize\textcolor{NordBlue}{\bfseries\insertsubsection}\hfill\null\newline%
    % \vskip3pt\null\hfill\noindent\textcolor{NordBrightBlack}{\rule[0.618ex]{0.618\linewidth}{0.618pt}}\hfill\null
}



\setbeamertemplate{footline}{
    \hfill\resizebox{16pt}{!}{\includegraphics{smallest_logo_transparent_bg.png}}\hskip3pt\null
}



% \setbeamertemplate{mini frames}[box]
% {\usebeamercolor{palette primary}}
% \setbeamertemplate{sidebar canvas}[vertical shading][top=primary.bg,middle=white,bottom=primary.bg]

\mdfdefinestyle{code}{%
	innertopmargin=0pt,
	innerbottommargin=12pt,
	skipabove=\baselineskip,
	skipbelow=\baselineskip,
	leftmargin=6pt,
	rightmargin=6pt,
	innerleftmargin=3,
	innerrightmargin=2,
	rightline=true,
	frametitlerule=true,
	backgroundcolor=NordDarkBlack,
	frametitlerulewidth=0pt,
	align=center,
	linecolor=NordBrightBlack,
	linewidth=1pt,
	shadow=false,
	shadowsize=3pt,
	shadowcolor=NordWhite,
	roundcorner=20pt
}

\setbeamersize{text margin left=2mm,text margin right=0mm}
\begin{document}

\setbeamertemplate{background canvas}{%
  \includegraphics[width = \paperwidth, height = \paperheight]
  {bg_gradient.png}%
}

\usemintedstyle{nord}


% \setminted[python]{texcomments=true,samepage=true,resetmargins=true,python3=true,tabsize=4,obeytabs=false,linenos=true,highlightcolor=NordBrightBlack,highlightlines={2, 5-9},fontsize=\tiny,autogobble,breaklines,breaksymbolleft=\textcolor{NordBrightBlack}{\tiny$\hookrightarrow$}}
\setminted[python]{
	texcomments=true,
	samepage=true,
	resetmargins=true,
	python3=true,
	tabsize=4,
	obeytabs=false,
	linenos=false,
	fontsize=\scriptsize,
	autogobble,
	breaksymbolleft=\textcolor{NordBrightBlack}{\tiny$\hookrightarrow$}
}

\setminted[bash]{
	texcomments=true,
	samepage=true,
	resetmargins=true,
	tabsize=4,
	obeytabs=false,
	linenos=false,
	fontsize=\normalsize,
	autogobble,
	breaksymbolleft=\textcolor{NordBrightBlack}{\tiny$\hookrightarrow$}
}

\section{Testing size}
\subsection[]{test file}
\begin{frame}[fragile=singleslide,noframenumbering,squeeze]
% \begin{frame}[noframenumbering]
    \frametitle{}
	\null\vfill
	\begin{center}
        \inputminted{python}{test.py}
	\end{center}
	\vfill\null
\end{frame}

% \section{Gerar Número Aleatório Entre 1 e 10}
% \subsection{\ttfamily\$RANDOM}
% \begin{frame}[noframenumbering]
% 	\null\vfill
% 	\begin{center}
%         \inputminted{bash}{random_number_1_10.sh}
% 	\end{center}
% 	\vfill\null
% \end{frame}

% \section{Mostrar primeiros e últimos elementos}
% \subsection{\ttfamily df.iloc}
% \begin{frame}[noframenumbering]
% 	\null\vfill
% 	\begin{center}
%         \inputminted{python}{df_head_tail.py}
% 	\end{center}
% 	\vfill\null
% \end{frame}

	% \vspace*{-3px}
	% \begin{tikzpicture}[%
	% 		overlay,%
	% 		remember picture,%
	% 		set style={{help lines}+=[very thin,NordBG]},%
	% 		pencildraw/.style={NordWhite,decorate,decoration={random steps,segment length=1pt,amplitude=0.15pt}},%
	% 		penciline={jag ratio=0.1}
	% ]
	% 	\draw[help lines] (0,0) grid (8,9);
	% 	\draw [<->, rounded corners, thick, NordMagenta] (0,2) -- (0,0) -- (3,0);
	% 	\draw[very thick, NordWhite] (0,0) to [out=90,in=195] (2,1.5);
	% 	\draw [<->,thick, NordCyan] (0,0) to [out=90,in=180] (1,1)
	% 	to [out=0,in=180] (2.5,0) to [out=0,in=-135] (4,1) ;
	% 	\draw[NordRed, thick] (0,9) -- (8,0);
	% 	\draw[NordBrightBlue, thick] (0,0) -- (8,9);
	% 	\filldraw[NordBG] (4,4.5) circle (1pt);
	% 	\node [above] at (4,4.5) {above};
	% 	\node [below] at (4,4.5) {below};
	% 	\node [left] at (4,4.5) {left};
	% 	\node [right] at (4,4.5) {right};

	% 	\filldraw[NordBG] (4,7.5) circle (1pt);
	% 	\node [below right] at (4,7.5) {below right};
	% 	\node [below left] at (4,7.5) {below left};
	% 	\node [above right] at (4,7.5) {above right};
	% 	\node [above left] at (4,7.5) {above left};

	% 	\draw [thick]
	% 	(2,9) -- (8,9);
	% 	\draw (2,8.8) -- (2, 9.2);
	% 	\draw (4,8.8) -- (4, 9.2);
	% 	\draw (6,8.8) -- (6, 9.2);
	% 	\draw (8,8.8) -- (8, 9.2);
	% 	\begin{scope}[NordBlack]
	% 		\node[align=left, below,fill=NordWhite,rounded corners] at (3,8.9)%
	% 		{This happens\\in period 1\\and is aligned\\ left};
	% 		\node[align=center, below,fill=NordWhite,rounded corners] at (5,8.9)%
	% 		{This happens\\in period 2\\and is centered};
	% 		\node[align=right, below,fill=NordWhite,rounded corners] at (7,8.9)%
	% 		{This happens\\in period 3\\and is right aligned};
	% 	\end{scope}

	% 	\draw[x=10ex,y=10ex] (6.5,6.5) sin (8,7.5);

	% 	\draw[snake=brace,thick] (0,3.73) -- (0.55,3.73);
	% 	\node[align=left,above right,NordBlack,fill=NordWhite,rounded corners] at (0.32,3.93)%
	% 	{Essa palavra chave indica a definição de uma classe};

	% 	\draw[->,pencildraw,thick,very nearly opaque] (0.61,2.53) to [out=80,in=180] (0.87,2.65);
	% 	\node[align=left,right,NordBlack,fill=NordWhite,rounded corners,inner sep=0.2ex] at (0.87,2.65)%
	% 	{\includegraphics[height=0.027\textheight]{def_func.png}};
	% 	% {Essa palavra chave indica a definição de uma função};

	% 	% \filldraw[NordWhite] (0.61,2.53) circle (1pt);
	% 	% \filldraw[NordWhite] (0.91,2.6) circle (1pt);
	% \end{tikzpicture}


    % \vskip12pt\hfill\resizebox{16px}{!}{\includegraphics{smallest_logo_transparent_bg.png}}\hskip4pt\null


% \section{Exemplo de Seção 1}
% \subsection{Exemplo de SubSeção 1}
% \begin{frame}[fragile=singleslide,noframenumbering,squeeze]
% 	\inputminted{python}{example.py}
% 	\vspace*{-3px}
% 	\begin{tikzpicture}[%
% 			overlay,%
% 			remember picture,%
% 			set style={{help lines}+=[very thin,NordBG]},%
% 			pencildraw/.style={NordWhite,decorate,decoration={random steps,segment length=1pt,amplitude=0.15pt}},%
% 			penciline={jag ratio=0.1}
% 	]
% 		\draw[help lines] (0,0) grid (8,9);
% 		\draw [<->, rounded corners, thick, NordMagenta] (0,2) -- (0,0) -- (3,0);
% 		\draw[very thick, NordWhite] (0,0) to [out=90,in=195] (2,1.5);
% 		\draw [<->,thick, NordCyan] (0,0) to [out=90,in=180] (1,1)
% 		to [out=0,in=180] (2.5,0) to [out=0,in=-135] (4,1) ;
% 		\draw[NordRed, thick] (0,9) -- (8,0);
% 		\draw[NordBrightBlue, thick] (0,0) -- (8,9);
% 		\filldraw[NordBG] (4,4.5) circle (1pt);
% 		\node [above] at (4,4.5) {above};
% 		\node [below] at (4,4.5) {below};
% 		\node [left] at (4,4.5) {left};
% 		\node [right] at (4,4.5) {right};

% 		\filldraw[NordBG] (4,7.5) circle (1pt);
% 		\node [below right] at (4,7.5) {below right};
% 		\node [below left] at (4,7.5) {below left};
% 		\node [above right] at (4,7.5) {above right};
% 		\node [above left] at (4,7.5) {above left};

% 		\draw [thick]
% 		(2,9) -- (8,9);
% 		\draw (2,8.8) -- (2, 9.2);
% 		\draw (4,8.8) -- (4, 9.2);
% 		\draw (6,8.8) -- (6, 9.2);
% 		\draw (8,8.8) -- (8, 9.2);
% 		\begin{scope}[NordBlack]
% 			\node[align=left, below,fill=NordWhite,rounded corners] at (3,8.9)%
% 			{This happens\\in period 1\\and is aligned\\ left};
% 			\node[align=center, below,fill=NordWhite,rounded corners] at (5,8.9)%
% 			{This happens\\in period 2\\and is centered};
% 			\node[align=right, below,fill=NordWhite,rounded corners] at (7,8.9)%
% 			{This happens\\in period 3\\and is right aligned};
% 		\end{scope}

% 		\draw[x=10ex,y=10ex] (6.5,6.5) sin (8,7.5);

% 		\draw[snake=brace,thick] (0,3.73) -- (0.55,3.73);
% 		\node[align=left,above right,NordBlack,fill=NordWhite,rounded corners] at (0.32,3.93)%
% 		{Essa palavra chave indica a definição de uma classe};

% 		\draw[->,pencildraw,thick,very nearly opaque] (0.61,2.53) to [out=80,in=180] (0.87,2.65);
% 		\node[align=left,right,NordBlack,fill=NordWhite,rounded corners,inner sep=0.2ex] at (0.87,2.65)%
% 		{\includegraphics[height=0.027\textheight]{def_func.png}};
% 		% {Essa palavra chave indica a definição de uma função};

% 		% \filldraw[NordWhite] (0.61,2.53) circle (1pt);
% 		% \filldraw[NordWhite] (0.91,2.6) circle (1pt);
% 	\end{tikzpicture}


% \end{frame}


% \begin{frame}{Font size in Beamer (default)}

% 	\begin{tikzpicture}[
% 	pencildraw/.style={
% 			white!75,
% 			decorate,
% 			decoration={random steps,segment length=0.8pt,amplitude=0.15pt}
% 		}
% 		]
% 		\node[pencildraw,draw] {Node};
% 		\draw[pencildraw] (0,1) -- (1,1) arc (0:-180:9mm and 7mm);
% 	\end{tikzpicture}

% 	\begin{tikzpicture}[penciline={jag ratio=1}]
% 		\draw[penciline={jag ratio=5},->](-5,0) -- (5,0);
% 		\draw[penciline={jag ratio=5},->](0,-5) -- (0,5);

% 		\draw[decorate,fill=blue!15,fill opacity=.25]  (-3,-3) rectangle (3,3);
% 		\draw[penciline={jag ratio=0},fill=green!15,fill opacity=.25] circle [radius=3cm];
% 		\draw[decorate,red,shift={(1,1)},rotate=45,fill=red!25,fill opacity=.5] %
% 		(1,1) -- (1,-1) -- (-1,-1) -- (-1,1) -- (1,1);
% 		\draw[penciline={jag ratio=0.01},color=blue]   plot (\x,{sin(\x r)}) node [decorate,draw] {$\sin(x)$} ;
% 	\end{tikzpicture}
% This is default font size

% \tiny This is tiny font size

% \scriptsize This is scriptsize font size

% \footnotesize This is footnotesize font size

% \small This is small font size

% \normalsize This is normalsize font size

% \large This is large font size

% \Large This is Large font size

% \LARGE This is LARGE font size

% \huge This is huge font size

% \Huge This is Huge font size

% \end{frame}

% \begin{frame}{Slanted and small caps text}
% This is \textsc{small caps text} and this is
% \textsl{slanted text}.\\~\\
% You can combine them, to produce \textsl{\textsc{small
% caps slanted text}} but also \textsc{\textbf{bold small caps}} or \textsl{\underline{underlined slanted text}}.\\~\\


% \emph{This} is emphasized and \textit{\emph{this} is
% also emphasized, although in a different way.}\\~\\


% Let $\bm{u}$, $\bm{v}$ be vectors and $\bm{A}$ be a
% matrix such that $\bm{Au}=\bm{v}$.
% This is a bold integral:
% \[
% \bm{\int_{-\infty}^{\infty} e^{-x^2}\,dx=\sqrt{\pi} }
% \]
% \end{frame}

% \begin{frame}{Text decorations provided by the \texttt{ulem} package}
% \uline{Underlined that breaks at the end of lines if they are too too long because the author won’t stop writing.} \\~\\
% \uuline{Double-Underlined text} \\~\\
% \uwave{Wavy-Underlined text} \\~\\
% \sout{Strikethrough text} \\~\\
% \xout{Struck with Hatching text} \\~\\
% \dashuline{Dashed Underline text} \\~\\
% \dotuline{Dotted Underline text}
% \end{frame}

% \begin{frame}{Basic Blocks}
% 	\begin{block}{Standard Block}
% 		This is a standard block.
% 	\end{block}
% 	\begin{alertblock}{Alert Message}
% 		This block presents alert message.
% 	\end{alertblock}
% 	\begin{exampleblock}{An example of typesetting tool}
% 		Example: MS Word, \LaTeX{}
% 	\end{exampleblock}
% \end{frame}
% % Frame 2
% \begin{frame}{Mathematical Environment Blocks}
% 	\begin{definition}
% 		This is a definition.
% 	\end{definition}

% 	\begin{theorem}
% 		This is a theorem.
% 	\end{theorem}

% 	\begin{lemma}
% 		This is a proof idea.
% 	\end{lemma}
% \end{frame}
% % Frame 3
% \begin{frame}{Mathematical Environment Blocks-Continued}
% 	\begin{proof}
% 		This is a proof.
% 	\end{proof}

% 	\begin{corollary}
% 		This is a corollary
% 	\end{corollary}

% 	\begin{example}
% 		This is an example
% 	\end{example}
% \end{frame}




% \begin{frame}{Vertical line between columns}

% 	\begin{columns}[T]%Tcb
% 		% Column 1
% 		\begin{column}{0.49\textwidth}

% 			\begin{itemize}
% 				\item Input layer: 2 neurons.
% 				\item Hidden layer: 5 neurons.
% 				\item Output layer: 2 neurons.
% 			\end{itemize}

% 		\end{column}
% 		% Column 2 (vertical line)
% 		\begin{column}{.02\textwidth}
% 			\rule{.1mm}{0.7\textheight}
% 		\end{column}
% 		% Column 3
% 		\begin{column}{0.49\textwidth}
% 			\includegraphics[width=\textwidth]{logo_borders.png}
% 		\end{column}

% 	\end{columns}
% \end{frame}

% \begin{frame}[plain,squeeze]{Defined Colors}{This is a subtitle}
% 	\begin{description}[Snow Storm]
% 		\item[Polar Night]
% 			\textcolor{NordDarkBlack}{NordDarkBlack} \quad \textcolor{NordBlack}{NordBlack}\\
% 			\textcolor{NordMediumBlack}{NordMediumBlack} \quad \textcolor{NordBrightBlack}{NordBrightBlack}
% 		\item[Snow Storm]
% 			\textcolor{NordWhite}{NordWhite} \quad \textcolor{NordBrighterWhite}{NordBrightestWhite}\\
% 			\textcolor{NordBrightestWhite}{NordBrightestWhite}
% 		\item[Frost]
% 			\textcolor{NordCyan}{NordCyan} \quad \textcolor{NordBrightCyan}{NordBrightCyan}\\
% 			\textcolor{NordBlue}{NordBlue} \quad \textcolor{NordBrightBlue}{NordBrightBlue}
% 		\item[Aurora]
% 			\textcolor{NordRed}{NordRed} \quad \textcolor{NordOrange}{NordOrange} \\
% 			\textcolor{NordYellow}{NordYellow} \quad \textcolor{NordGreen}{NordGreen} \\
% 			\textcolor{NordMagenta}{NordMagenta}
% 	\end{description}
% \end{frame}

% \begin{frame}
% \begin{tikzpicture}
% \draw[NordRed, thick] (-1,2) -- (2,-4);
% \draw[NordBrightBlue, thick] (-1,-1) -- (2,2);
% \filldraw[black] (0,0) circle (2pt) node[anchor=west]{
% 		\includegraphics[width=0.5cm]{logo_borders.png}
% };
% \end{tikzpicture}
% 	% \begin{tikzpicture}[overlay,remember picture]
% 	% \node[right=0.2cm] at (current page.150) {%
% 	% 	\includegraphics[width=0.5cm]{logo_borders.png}
% 	% };
% 	% \end{tikzpicture}
% \end{frame}

% \begin{frame}[fragile]
% \tiny\ttfamily
% \begin{lstlisting}
% #!/bin/bash
% echo the $# parameter did not destroy pygments syntax highlighting
% \end{lstlisting}
% \end{frame}

% \lstset{basicstyle=\footnotesize\ttfamily,breaklines=true}
% \begin{frame}
% 	\lstinputlisting[language=Bash]{bash_script.sh}
% \end{frame}


% \begin{frame}

%     This is your first presentation!
% 	\newline
%     \rmfamily This is your first presentation!
% 	\newline
% 	\textup{Upright} \textit{Italic} \textsl{Slanted}
% 	\newline


%     \sffamily This is your first presentation!
% 	\newline
% 	\textsf{Upright  \textit{Italic} \textsl{Slanted}}
% 	\newline

%     \ttfamily This is your first presentation!
% 	\newline
% 	\texttt{Upright  \textit{Italic} \textsl{Slanted}}
% 	\newline

% 	\textsc{Small Capitals}

% \end{frame}

% % \begin{overlayarea}<overlay spec>{area width}{area height}
% \begin{overlayarea}{1080px}{1080px}
%   content
% \end{overlayarea}

\end{document}
